\section{Discussion}
\label{sec:discussion}

\paragraph{What the signature establishes.}
A signature summarizes trajectory--correctness associations under a fixed
model-family, prompt, evaluator, and extraction protocol. Matched Qwen
checkpoints differ under that protocol, but budget and scale audits show that
fine nearest-neighbor topology is neither a universal fingerprint nor a causal
taxonomy.

\paragraph{Association is not mechanism or utility.}
Geometry may reflect planning, formatting, length, termination, confidence, or
unmeasured variables; our controls remove specified confounds, not all causal
alternatives. Association requires a reproducible correct--incorrect difference;
selection requires a same-pool gain over agreement, length, and likelihood;
intervention requires a locked paired gain over fair controls without behavioral
degradation. Evidence at one level does not establish the next.

\paragraph{Implications for transfer.}
Source steering is useful only if it competes with target calibration and
standard activation steering. Under matched target, prompt, split, and budget,
source \textsc{CrossSteer} does not do so. The five-label target-local point is
significant after family-wise correction over the six predeclared label budgets,
but the non-monotone curve is a bounded cell-level result, not a paper-wide
intervention claim or a general sample-efficiency law.

\paragraph{Conclusion.}
Post-training variants can exhibit different length-residualized trajectory
signatures without yielding a useful selector or transferable intervention;
the contribution is a measurement object and an evidence hierarchy that
keeps observability, selection utility, and intervention efficacy distinct.
